\begin{textsample}{POS Dim 1 – vem\_ed\_17 – Score 10.00 – vem\_ed\_17/t074.txt}  \label{ex:f1_pos_133}
QUEM \textbf{É} QUEM Hope Windle \textbf{é} diretora do SUNY COIL Center (State University of New York / EUA), \textbf{referência} mundial em Intercâmbios Virtuais.
A \textbf{seguir}, tem-se uma síntese da entrevista concedida por videoconferência para VEm em 23 de março de 2023.
Sabemos \textbf{que} o SUNY COIL Center está de \textbf{casa} nova, SUNY Oneonta.
Pode nos falar sobre a \textbf{importância} disso?
Agora estamos conectados com estudantes e professores, ao redor do mundo, nacionalmente e dentro da SUNY.
Temos o \textbf{que} chamamos “colaboratório”, algo como uma comunidade de prática, ou incubadora. \textbf{É} um caminho para \textbf{compartilhar} ideias em nossa comunidade de 153 instituições ao redor do mundo.
Também queremos atrair os alunos quando fazemos apresentações ou eventos, para \textbf{que} tenhamos \textbf{pessoas} localmente também, porque eu estou sempre no Zoom, virtualmente.
Qual \textbf{foi} um dos projetos mais desafiadores com \textbf{que} você trabalhou?
Essa \textbf{é} uma boa \textbf{questão}.
Houve alguns, eu diria \textbf{que} quando \textbf{fui} coordenadora na SUNY Ulster [ 2017-2022 ], fizemos um projeto com um adorável professor de Belarus, ele queria abordar empreendedorismo e turismo.
Mas, durante a colaboração, muitos dos links \textbf{compartilhados} com os alunos \textbf{eram} censurados pelo governo (e não \textbf{sabíamos} disso).
Não \textbf{era} \textbf{permitido} aos estudantes \textbf{compartilhar} vídeos do YouTube e o professor não podia realmente \textbf{contar} o \textbf{que} estava acontecendo.
Os estudantes comentaram: \textbf{é} uma grande ideia empreender projetos turísticos reais, mas não podemos fazer neste país. \textbf{Foi} interessante \textbf{reconhecer} diferentes \textbf{formas} de governo e como afetam os negócios.
Como a sua experiência de designer instrucional influencia a \textbf{forma} como você dirige o SUNY COIL Center?
Quando se pensa em liderança no COIL, temos Jon Rubin, \textbf{que} \textbf{foi} professor e veio com um \textbf{ponto} de vista de professor: pensar no \textbf{conteúdo}, como juntar os professores e desenvolver os projetos.
Hope Windle, diretora do SUNY COIL Center Quando Mary Lou Forward chegou no SUNY COIL Center [ 2017 ], abrimos para um \textbf{número} muito maior de membros.
E então, como resultado da pandemia para mim, minha \textbf{vivência} em design instrucional me faz querer \textbf{que} a experiência COIL \textbf{seja} flexível para \textbf{todos}, em termos das diferentes \textbf{formas} de fazer, além de fornecer o kit de ferramentas e os recursos para \textbf{que} as \textbf{pessoas} possam visualizar facilmente.
O \textbf{que} você prevê para os projetos COIL nos próximos cinco anos?
Estamos atingindo um \textbf{momento} incrível.
No começo dos projetos COIL [ 2006 ], a pergunta \textbf{era}: como fazemos, como usamos a tecnologia?
Agora, temos \textbf{questões} ligadas à esperança \textbf{crítica}, para \textbf{que} \textbf{pessoas}, especialmente aquelas com dificuldades econômicas, possam redesenhar sua educação.
COIL \textbf{é} um meio para \textbf{que} repensem o \textbf{que} vão fazer de suas vidas.

% matched lemmas: casa, compartilhar, contar, conteúdo, crítico, forma, importância, momento, número, permitir, pessoa, ponto, que, questão, reconhecer, referência, saber, seguir, ser, todo, vivência
\end{textsample}
